% ---------------------------------------------------------------------------
% Preambolo condiviso della monografia NanoFaaS.
% Compilazione: latexmk -pdf nanofaas.tex   (pdfLaTeX, TeX Live 2026+)
% ---------------------------------------------------------------------------

\usepackage[T1]{fontenc}
\usepackage[utf8]{inputenc}
\usepackage[italian]{babel}
\usepackage{lmodern}
\usepackage{microtype}

\usepackage{geometry}
\geometry{a4paper, top=3cm, bottom=3cm, inner=3cm, outer=2.5cm}

\usepackage{amsmath, amssymb}
\usepackage{booktabs}
\usepackage{longtable}
\usepackage{array}
\usepackage{multirow}
\usepackage{tabularx}
\usepackage{caption}
\usepackage{subcaption}
\usepackage{enumitem}
\usepackage{float}
\usepackage{adjustbox}

% Figure TikZ: le sorgenti in figures/ sono documenti `standalone`.
% Caricando il pacchetto `standalone` qui, \input{figures/fig-x} le include
% ignorandone il \documentclass e riusando questo preambolo.
\usepackage{figures/nanofaas-figs}
\usepackage{standalone}
\usepackage{pgfplots}
\pgfplotsset{compat=1.18}

\usepackage{listings}
\usepackage{xcolor}

\usepackage[hidelinks]{hyperref}
\usepackage[nameinlink,italian]{cleveref}

% ---- stile dei listati -----------------------------------------------------
\lstdefinestyle{nf}{
  basicstyle=\ttfamily\footnotesize,
  keywordstyle=\color{nfblue}\bfseries,
  commentstyle=\color{nfslate}\itshape,
  stringstyle=\color{nfgreen},
  numbers=left, numberstyle=\tiny\color{nfslate}, numbersep=7pt,
  frame=leftline, framerule=1pt, rulecolor=\color{nfslate!40},
  xleftmargin=10pt, breaklines=true, breakatwhitespace=false,
  postbreak=\mbox{\textcolor{nfslate}{$\hookrightarrow$}\space}, columns=flexible,
  showstringspaces=false, tabsize=2, captionpos=b,
  literate={à}{{\`a}}1 {è}{{\`e}}1 {é}{{\'e}}1 {ì}{{\`i}}1 {ò}{{\`o}}1 {ù}{{\`u}}1
           {—}{{---}}1 {–}{{--}}1 {’}{{'}}1 {‘}{{'}}1 {“}{{"}}1 {”}{{"}}1 {≥}{{$\geq$}}1 {≤}{{$\leq$}}1 {×}{{$\times$}}1 {λ}{{$\lambda$}}1 {μ}{{$\mu$}}1,
}
\lstset{style=nf}
\lstdefinelanguage{yamlnf}{
  keywords={true,false,null,enabled,mode,strategy,metric},
  sensitive=false, comment=[l]{\#}, morestring=[b]',morestring=[b]"
}

% ---- elementi ricorrenti ---------------------------------------------------
\newcommand{\code}[1]{\texttt{#1}}
\newcommand{\nf}{NanoFaaS}

% Riquadro "da misurare": marca un dato non ancora rilevato.
\usepackage[most]{tcolorbox}
\newtcolorbox{damisurare}[1][]{
  colback=nfamberbg, colframe=nfamber, boxrule=0.6pt, arc=2pt,
  fonttitle=\bfseries\sffamily\small, title={Da misurare}, #1}
\newcommand{\todomis}[1]{\textcolor{nfamber}{\textbf{[da misurare: #1]}}}

% Numerazione profonda ma indice compatto.
% Tolleranza extra: il testo italiano con molti \code{...} non spezzabili
% produrrebbe altrimenti righe che escono dal margine.
\setlength{\emergencystretch}{6em}
\hyphenpenalty=1000
\tolerance=3000

\setcounter{secnumdepth}{3}
\setcounter{tocdepth}{1}
